\chapter{Carry à travers les classes d'actifs}\label{ch:s1:carry-across-asset-classes}

Emprunter dans une devise qui rapporte peu, prêter dans une devise qui rapporte beaucoup, et encaisser la différence tant que le taux de change ne
joue pas contre le porteur. En général il ne bouge pas assez, et le carry trade de change rapporte depuis des décennies~; puis, en quelques semaines, les devises à haut rendement chutent
ensemble, et des années de revenus s'envolent. Koijen, Moskowitz, Pedersen et Vrugt ont montré que la même idée, le rendement qu'un actif procure si son prix ne change pas,
prédit les rendements des actions, des obligations, des matières premières, du crédit et des options aussi bien que des devises, et que les \omterm{def:s1:carry-across-asset-classes:carry}{stratégies de portage} de toutes les classes se comportent mal ensemble
en période de récession mondiale. Sur l'univers synthétique, un portefeuille de \omterm{def:s1:carry-across-asset-classes:diversified}{carry diversifié} dégage un ratio de Sharpe de 0.95 lorsque tout le carry est encaissé, et perd 26~\% de son
capital, à 10~\% de volatilité, lors du deuxième krach boursier. La construction est \texttt{firm.carrystrat}.

\section{Le carry partout}

Le livre~1, chapitre~7 définit le carry comme ce qu'une position rapporte si les prix restent où ils sont. Pour une position sur contrats à terme, il prend une forme commune~: avec un prix au comptant ou
de référence $S$ et un prix à terme $F$ pour une livraison dans $\tau$ années, le carry vaut $C = (S - F)/(F\tau)$, soit la convergence du prix à terme vers le prix au comptant
si le prix au comptant ne bouge pas. Chaque classe lui donne un nom.
\begin{itemize}
\item \textbf{Devises}~: le déport à terme, soit le taux d'intérêt de la devise étrangère moins celui de la devise domestique.
\item \textbf{Indices actions}~: le rendement de dividende moins le taux de financement.
\item \textbf{Obligations d'État}~: le rendement au-dessus du taux de financement, plus le roll-down, soit le gain de prix d'une obligation qui vieillit en descendant une courbe ascendante.
\item \textbf{Matières premières}~: le rendement de roll, positif en backwardation (prix à terme inférieur au prix au comptant) et négatif en contango, fixé par les coûts de stockage et le
rendement d'opportunité (convenience yield) (livre~3, chapitre~10).
\end{itemize}
Koijen et ses coauteurs écrivent le rendement en excès espéré d'une position à terme comme son carry plus la variation espérée du prix au comptant. Si les prix au comptant ne
bougent pas en moyenne, tout le carry est encaissé~; s'ils bougeaient pour l'annuler, comme le supposeraient l'hypothèse des anticipations et la parité des taux d'intérêt non couverte,
aucun ne le serait. \texttt{firm.synthfut} introduit dans le rendement espéré une part $\lambda$ du carry~: $\lambda = 0.5$ dans l'univers de tous ces chapitres, et
$\lambda = 1$ dans une variante où les prix ne jouent pas contre le carry.

\begin{definition}[Stratégie de portage]\label{def:s1:carry-across-asset-classes:carry}
Une \emph{stratégie de portage}\index{stratégie de portage} (carry strategy) détient des positions longues sur les actifs d'une classe au carry le plus élevé et des positions courtes sur ceux qui ont
le carry le plus faible, avec des pondérations croissant avec le rang du carry, rééquilibrées à mesure que les niveaux de carry changent~; elle encaisse le spread de carry et supporte le risque que les prix
évoluent à son encontre.
\end{definition}

Dans l'univers synthétique, le carry se lit sur la courbe (\cref{lst:s1:carry-across-asset-classes:measures})~: la pente logarithmique entre les deux premiers
contrats mensuels séparés de 21 jours de négociation et annualisée. Pour la plupart des marchés, cela égale exactement le carry introduit, car les courbes du modèle sont des droites en
maturité. Pour les trois matières premières saisonnières, ce n'est pas le cas~: leurs courbes portent le cycle annuel de leurs prix au comptant, et la pente à un mois indique surtout
la direction que prend la saison. Sa corrélation avec le vrai carry vaut 0.05. La pente entre contrats espacés de douze mois annule le cycle et
a une corrélation de 1.00. Les portefeuilles de ce chapitre l'utilisent pour les matières premières.

\section{Portefeuilles de carry par classe}

Le portefeuille de chaque classe pondère ses dix marchés par le rang centré de leur carry, avec une exposition brute de 1, rééquilibré chaque jour à 2, 1, 2 et 4 points
de base par unité échangée, selon la classe (\cref{lst:s1:carry-across-asset-classes:books})~; chacun est ramené sur l'ensemble de l'échantillon à une volatilité de 10~\%, et le
portefeuille diversifié attribue à chaque classe le même poids.

\begin{definition}[Carry diversifié]\label{def:s1:carry-across-asset-classes:diversified}
Le \emph{carry diversifié}\index{carry diversifié} (diversified carry) est un portefeuille de \omterm{def:s1:carry-across-asset-classes:carry}{stratégies de portage} dans plusieurs classes d'actifs, chacune ramenée à un risque comparable, qui
encaisse les primes de carry de toutes les classes à une volatilité inférieure à celle de chacune, tant que leurs rendements ne sont pas corrélés.
\end{definition}

\begin{center}
\small
\begin{tabular}{@{}lccccc@{}}
\toprule
Ratio de Sharpe après coûts & actions & obligations & devises & matières premières & diversifié\\
\midrule
moitié du carry encaissé ($\lambda = 0.5$) & 0.17 & 0.22 & 0.03 & 0.27 & 0.34\\
tout le carry encaissé ($\lambda = 1$) & 0.28 & 0.53 & 0.49 & 0.59 & 0.95\\
\midrule
asymétrie mensuelle ($\lambda = 1$) & 0.13 & $-0.09$ & $-0.36$ & $-0.16$ & $-0.20$\\
\bottomrule
\end{tabular}
\end{center}

Deux éléments fixent les chiffres. La part de carry encaissé double ou triple le ratio de Sharpe de chaque classe (\cref{fig:s1:carry-across-asset-classes:classes}),
c'est pourquoi la question empirique de savoir si les prix jouent contre le carry est toute la question. Et le portefeuille diversifié bat chaque classe~: quatre portefeuilles
peu corrélés en temps normal, chacun à 10~\% de volatilité, se combinent en un portefeuille dont le ratio de Sharpe est proche du double de celui de la classe moyenne. Le
portefeuille de devises à $\lambda = 0.5$ rapporte presque rien, car ses krachs, décrits plus loin, reprennent ce que rapporte son carry.

\begin{omfigure}
\begin{tikzpicture}
\begin{axis}[width=10.5cm,height=5.2cm,ybar,bar width=11pt,ylabel={\small ratio de Sharpe après coûts},tick label style={font=\footnotesize},
  symbolic x coords={equity,bond,currency,commodity,diversified},xtick=data,xticklabels={actions,obligations,devises,{matières\\premières},diversifié},x tick label style={align=center},ymin=0,ymax=1.1,enlarge x limits=0.14,grid=major,
  grid style={gray!25},legend columns=2,legend style={font=\scriptsize,at={(0.5,-0.3)},anchor=north,draw=none,fill=none,
  /tikz/every even column/.append style={column sep=8pt}},area legend,table/col sep=comma]
\addplot[fill=omExo!50,draw=omExo] table[x=book,y=half]{figdata/strategies-1/20-carry-across-asset-classes/classes.csv};
\addplot[fill=omThm!60,draw=omThm] table[x=book,y=full]{figdata/strategies-1/20-carry-across-asset-classes/classes.csv};
\legend{moitié du carry encaissé,tout le carry encaissé}
\end{axis}
\end{tikzpicture}

\omcaption{Portefeuilles de carry sur l'univers synthétique, trente ans~: ratios de Sharpe après coûts par classe et pour le portefeuille diversifié, lorsque la moitié du carry de chaque
marché est encaissée comme rendement espéré, et lorsque tout l'est. Données~: \texttt{s1\_carry.summary}.}\label{fig:s1:carry-across-asset-classes:classes}
\end{omfigure}

\section{Krachs du carry}

\begin{definition}[Krach du carry]\label{def:s1:carry-across-asset-classes:crash}
Un \emph{krach du carry}\index{krach du carry} (carry crash) est une perte soudaine et importante d'une \omterm{def:s1:carry-across-asset-classes:carry}{stratégie de portage} lorsque les actifs à fort carry chutent et que les actifs à faible carry montent ensemble,
généralement lorsque les positions de carry à effet de levier sont débouclées lors d'un recul de l'appétit pour le risque ou de la liquidité de financement~; il en résulte que les rendements de la stratégie sont négativement asymétriques.
\end{definition}

Brunnermeier, Nagel et Pedersen ont documenté ce schéma sur les devises~: les variations de taux de change entre devises à taux d'intérêt élevé et devises à taux bas sont négativement
asymétriques, car les carry trades sont débouclés brutalement lorsque l'appétit pour le risque et la liquidité de financement reculent. Lustig, Roussanov et Verdelhan ont trouvé pourquoi la prime
existe~: les devises à taux élevé sont davantage exposées à un facteur de risque mondial, si bien que l'investisseur en carry porte un risque mondial, surtout dans les mauvais moments. La prime est
une compensation pour les pertes survenant quand elles font le plus mal.

\begin{dated}{2026-09}{Le débouclage du carry d'août 2024}\label{dat:s1:carry-across-asset-classes:yen}
La Banque des règlements internationaux a décrit la turbulence des marchés du début d'août 2024 dans sa Quarterly Review de septembre 2024~: après les réunions de la Réserve
fédérale et de la Banque du Japon, jugées quelque peu restrictives, les marchés ont surréagi à une publication décevante sur le marché du travail américain, et les carry trades de devises
ont été débouclés au milieu de l'évolution des anticipations de taux et d'une volatilité accrue, et les devises de financement, principalement le yen, se sont fortement appréciées pendant un court moment.
La revue cite une estimation des positions de change à terme des hedge funds, d'environ 160~milliards de dollars.
\end{dated}

L'univers synthétique introduit le krach~: à chaque krach boursier, les devises à fort carry perdent 8~\% par unité de score z de leur carry sur les cent derniers jours.
Le portefeuille de devises perd 29.0~\%, 24.9~\% et 28.0~\% lors des trois krachs ($\lambda = 1$, à 10~\% de volatilité) et présente l'asymétrie mensuelle la plus négative.
Les krachs touchent aussi les autres classes, car ils font bouger les prix d'une manière qui recoupe le carry. Le portefeuille diversifié perd 12.9~\%, 26.3~\% et
13.1~\%~: lors du deuxième krach, il a perdu davantage que n'importe quelle classe prise isolément, car toutes les classes ont perdu en même temps (\cref{fig:s1:carry-across-asset-classes:paths}). Cela
est en miniature le constat de Koijen et de ses coauteurs~: les classes sont en moyenne décorrélées et deviennent corrélées au moment où cela compte.

\begin{omfigure}
\begin{tikzpicture}
\begin{axis}[width=11cm,height=5.4cm,xlabel={\small année},ylabel={\small rendement cumulé (\%)},tick label style={font=\footnotesize},
  xmin=1,xmax=30,ymin=-80,ymax=320,grid=major,grid style={gray!25},legend columns=2,legend style={font=\scriptsize,at={(0.5,-0.3)},
  anchor=north,draw=none,fill=none,/tikz/every even column/.append style={column sep=8pt}},table/col sep=comma]
\fill[omExo!15] (axis cs:7.5,-80) rectangle (axis cs:7.897,320);
\fill[omExo!15] (axis cs:18,-80) rectangle (axis cs:18.397,320);
\fill[omExo!15] (axis cs:27,-80) rectangle (axis cs:27.397,320);
\addplot[omThm,thick] table[x=year,y=diversified]{figdata/strategies-1/20-carry-across-asset-classes/paths.csv};
\addplot[omDef,thick] table[x=year,y=currency]{figdata/strategies-1/20-carry-across-asset-classes/paths.csv};
\legend{carry diversifié,carry de devises}
\end{axis}
\end{tikzpicture}

\omcaption{Les livres de \omterm{def:s1:carry-across-asset-classes:diversified}{carry diversifié} et de carry de change sur l'univers synthétique lorsque tout le carry est perçu, chacun à une volatilité de 10~\%, sous forme de sommes cumulées de
rendements quotidiens~; les bandes grisées sont les trois krachs boursiers. Données~: \texttt{s1\_carry.paths}.}\label{fig:s1:carry-across-asset-classes:paths}
\end{omfigure}

Le carry et le trend forment une paire naturelle. Les rendements quotidiens du livre de \omterm{def:s1:carry-across-asset-classes:diversified}{carry diversifié} présentent une corrélation de $-0.00$ avec le livre de trend mixte du chapitre~19
($\lambda = 0.5$), et lors des krachs ils évoluent en sens opposés~: le trend gagne, le carry perd. Un livre qui détient les deux a des rendements plus réguliers que
l'un ou l'autre pris isolément.

\section{Carry sur la courbe du WTI}

Les règlements-livraisons des deuxième et troisième contrats WTI, séparés d'un mois, fournis par l'EIA, donnent le carry du pétrole brut de 1985 à 2024 (la \cref{fig:s1:carry-across-asset-classes:wti}).
Utiliser ces contrats plutôt que le premier contrat permet d'éviter le contrat dans ses derniers jours avant l'échéance. La moyenne s'élève à 1.2~\% par an~; la courbe était en
backwardation pendant 51.7~\% des jours. Les valeurs extrêmes furent de $+126~\%$ par an le 14 janvier 1991 et de $-575~\%$ le 21 avril 2020, le lendemain du jour où le contrat de mai
a été réglé en territoire négatif (livre~3, chapitre~2). Une règle qui détient le contrat à terme reconduit en position longue en backwardation et en position courte en contango, avec une volatilité
cible de 40~\% et 2 points de base par transaction a rapporté 24.8~\% par an de 1986 à 2024, avec un ratio de Sharpe de 0.61~; il s'agit d'un seul marché, dans l'échantillon, sans paramètre à
ajuster au-delà du signe.

\begin{omfigure}
\begin{tikzpicture}
\begin{axis}[width=11cm,height=5cm,xlabel={\small année},ylabel={\small carry, contrats 2--3 (\% par an)},tick label style={font=\footnotesize},
  xmin=1985,xmax=2024.5,ymin=-200,ymax=150,grid=major,grid style={gray!25},table/col sep=comma,/pgf/number format/1000 sep={}]
\addplot[omThm] table[x=year,y=carry_pct]{figdata/strategies-1/20-carry-across-asset-classes/wti_carry.csv};
\addplot[black!40,dashed] coordinates {(1985,0) (2024.5,0)};
\end{axis}
\end{tikzpicture}

\omcaption{Le carry des contrats à terme sur pétrole brut WTI du NYMEX, $12\ln(F_2/F_3)$, à partir des règlements quotidiens de l'EIA pour les deuxième et troisième contrats, un jour de négociation sur
cinq, de 1985 à 2024~; positif en backwardation. L'échelle s'arrête à $-200~\%$~: le plus bas quotidien fut de $-575~\%$ le 21 avril 2020. Données~:
\texttt{s1\_carry.wti}.}\label{fig:s1:carry-across-asset-classes:wti}
\end{omfigure}

\section{Fiches de stratégie}

\begin{strategyfile}{Carry de change}\label{strat:s1:carry-across-asset-classes:fx}
\sfield{Qui vous paie, et pourquoi} Les investisseurs qui se couvrent ou se financent dans des devises à faibles taux, et le prix de marché du risque de krach global.
\sfield{Instruments et plateformes} Contrats de change à terme et contrats à terme sur devises développées et émergentes liquides.
\sfield{Signal} Le report de terme (écart de taux d'intérêt), classé.
\sfield{Dimensionnement et exécution} Position longue dans les devises à taux élevés, courte dans celles à taux faibles, pondérée par la volatilité~; termes mensuels reconduits.
\sfield{Coûts} Faibles~; points de terme et fourchettes.
\sfield{Comment elle meurt} Krachs lorsque les positions de carry à effet de levier se débouclent ensemble.
\sfield{Horizon, capacité, infrastructure} Mois~; grande capacité sur les principales devises.
\sfield{Backtest honnête} Taux de change à terme tels que cotés, et non des substituts de taux d'intérêt~; les années de krach de l'échantillon.
\sfield{Sources} Brunnermeier, Nagel et Pedersen (2008)~; Lustig, Roussanov et Verdelhan (2011)~; l'encadré daté.
\end{strategyfile}

\begin{strategyfile}{Carry obligataire}\label{strat:s1:carry-across-asset-classes:bond}
\sfield{Qui vous paie, et pourquoi} Les investisseurs qui acceptent une prime de terme pour détenir une duration courte.
\sfield{Instruments et plateformes} Futures sur obligations d'État de différents pays.
\sfield{Signal} Le rendement au-delà du taux de financement, plus le roll-down, classé.
\sfield{Dimensionnement et exécution} Long--short neutre en duration ou pondéré par la volatilité.
\sfield{Coûts} Faibles.
\sfield{Comment elle meurt} Hausses brutales des taux à l'avant des courbes pentues.
\sfield{Horizon, capacité, infrastructure} Mois~; un modèle de courbe par pays.
\sfield{Backtest honnête} Carry tiré de la courbe connue à l'époque, avec le moins cher à livrer de chaque contrat à terme.
\sfield{Sources} Koijen, Moskowitz, Pedersen et Vrugt (2018).
\end{strategyfile}

\begin{strategyfile}{Rendement de roll sur matières premières}\label{strat:s1:carry-across-asset-classes:commodity}
\sfield{Qui vous paie, et pourquoi} Les hedgers qui paient pour se décharger du risque d'inventaire~; le prix de marché de la rareté.
\sfield{Instruments et plateformes} Contrats à terme sur matières premières.
\sfield{Signal} La pente de la courbe, sur douze mois pour les matières premières saisonnières.
\sfield{Dimensionnement et exécution} Position longue sur les matières premières en backwardation, courte sur celles en contango, pondérée par la volatilité.
\sfield{Coûts} Modérés~; roll.
\sfield{Comment elle meurt} Excédents de stockage et squeezes qui inversent la relation entre pente et rendement.
\sfield{Horizon, capacité, infrastructure} Mois.
\sfield{Backtest honnête} Les contrats qui auraient été détenus, et non les contrats les plus proches~; pentes saisonnières séparées du carry.
\sfield{Sources} Koijen, Moskowitz, Pedersen et Vrugt (2018)~; le test saisonnier de ce chapitre.
\end{strategyfile}

\begin{strategyfile}{Carry diversifié}\label{strat:s1:carry-across-asset-classes:diversified}
\sfield{Qui vous paie, et pourquoi} Les primes de carry de toutes les classes à la fois.
\sfield{Instruments et plateformes} Contrats à terme et contrats de change à terme sur actions, obligations, devises et matières premières.
\sfield{Signal} Carry dans chaque classe, classé au sein de la classe.
\sfield{Dimensionnement et exécution} Chaque classe au même risque~; rééquilibrage mensuel ou dès que les carry changent.
\sfield{Coûts} Faibles à modérés.
\sfield{Comment elle meurt} Pertes corrélées lors des récessions mondiales, lorsque la diversification disparaît.
\sfield{Horizon, capacité, infrastructure} Mois~; un pipeline de données pour le carry de chaque classe.
\sfield{Backtest honnête} Mesures de carry définies avant le test~; récessions de l'échantillon.
\sfield{Sources} Koijen, Moskowitz, Pedersen et Vrugt (2018)~; ce chapitre~: 0.95 avec tout le carry perçu, une perte de 26.3~\% lors d'un krach.
\end{strategyfile}

\begin{strategyfile}{Carry de la courbe du WTI}\label{strat:s1:carry-across-asset-classes:wti}
\sfield{Qui vous paie, et pourquoi} Les producteurs et les consommateurs qui se couvrent, et qui paient pour l'immédiateté en backwardation et pour le stockage en contango.
\sfield{Instruments et plateformes} Contrats à terme WTI du NYMEX, reconduits.
\sfield{Signal} Le signe de la pente entre les deuxième et troisième contrats.
\sfield{Dimensionnement et exécution} Position longue en backwardation, courte en contango, à une cible de volatilité.
\sfield{Coûts} Roll et fourchettes.
\sfield{Comment elle meurt} Comme il s'agit d'un seul marché, elle n'a aucune diversification~; limites de stockage (avril 2020).
\sfield{Horizon, capacité, infrastructure} De quelques semaines à quelques mois.
\sfield{Backtest honnête} La série reconduite~; la pente issue de contrats éloignés de l'échéance.
\sfield{Sources} Règlements-livraisons de l'EIA~; ce chapitre~: un ratio de Sharpe de 0.61, de 1986 à 2024, dans l'échantillon.
\end{strategyfile}

\section{Tutoriel~: automne sur le marché synthétique}\label{tut:s1:carry-across-asset-classes:tut}

\begin{tutorial}
\textbf{Objectif.} Lire le carry sur les courbes synthétiques, construire des livres de carry par classe et diversifié selon deux hypothèses sur la part du carry perçue,
les mesurer lors des krachs, et lire le carry de la courbe de l'EIA pour le pétrole brut. \textbf{État final~:} le tableau et les trois figures.
\begin{tutsteps}
\item \textbf{Mesures de carry et pondérations par rang}.
\omcode{code/firm/carrystrat/firm_carrystrat.py}{26}{47}{Carry par classe et pondérations par rang au sein des groupes.}\label{lst:s1:carry-across-asset-classes:measures}
\item \textbf{Les livres}~: un par classe, puis le livre diversifié.
\omcode{code/strategies-1/20-carry-across-asset-classes/python/s1_carry.py}{60}{72}{Livres de carry par classe et diversifié.}\label{lst:s1:carry-across-asset-classes:books}
\item \textbf{Exécuter} \texttt{carry\_quality()}, \texttt{summary} avec le signal mixte aux primes 0.5 et 1.0, \texttt{trend\_correlation()},
\texttt{wti()}, et \texttt{fig\_carry.py}.
\end{tutsteps}
\textbf{Pour aller plus loin.} Dimensionner le livre de chaque classe ex ante avec une prévision EWMA plutôt que dans l'échantillon~; ajouter un filtre de krach qui coupe le carry de change
lorsque la volatilité bondit~; combiner le carry au livre de trend du chapitre~19 et mesurer les drawdowns combinés.
\end{tutorial}

\section{Construction~: stratégies de carry}\label{bld:s1:carry-across-asset-classes:carrystrat}

\begin{build}
\bfield{Rôle} Mesures de carry pour toutes les classes, livres de carry par rang et diagnostics de krach.
\bfield{Interface} \texttt{curve\_carry(f\_near, f\_far, days)}, \texttt{forward\_discount(log\_spot, log\_fwd, days)},
\texttt{dividend\_carry(div\_yield, rate)}, \texttt{rank\_weights(x, groups)}, \texttt{book(r, w, cost)}, \texttt{scale(pnl, target, start)}, \texttt{windows(pnl, spans)},
\texttt{skew\_monthly(pnl, start, days)}.
\bfield{Règles} Carry annualisé sur 252 jours~; rangs au sein de chaque groupe~; pondérations avec une exposition brute égale à un par groupe~; P\&L issu des pondérations fixées à la clôture précédente.
\bfield{Tests d'acceptation} \texttt{code/firm/carrystrat/tests/}~: carry à partir de deux cours, d'un forward et d'un rendement de dividende calculés à la main~; pondérations par rang au sein de chaque groupe~;
P\&L, coûts, fenêtres et série asymétrique.
\bfield{Pour aller plus loin} Dimensionnement ex ante~; carry à partir de courbes complètes avec roll-down~; filtres de krach.
\end{build}

\begin{omsources}
\item R.~S.~J.~Koijen, T.~J.~Moskowitz, L.~H.~Pedersen et E.~B.~Vrugt, «~Carry~», \emph{Journal of Financial Economics} 127(2), 2018.
\item M.~K.~Brunnermeier, S.~Nagel et L.~H.~Pedersen, «~Carry trades and currency crashes~», \emph{NBER Macroeconomics Annual} 23, 2008.
\item H.~Lustig, N.~Roussanov et A.~Verdelhan, «~Common risk factors in currency markets~», \emph{Review of Financial Studies} 24(11), 2011.
\item Bank for International Settlements, «~Carry off, carry on~», \emph{BIS Quarterly Review}, septembre 2024.
\end{omsources}

\section{Exercices}

\begin{exercise}[$\star$]\label{exo:s1:carry-across-asset-classes:1}
Deux contrats à terme espacés de 21 jours de négociation sont cotés 100 et 98. Calculer le carry annualisé sur 252 jours.
\end{exercise}

\begin{exercise}[$\star$]\label{exo:s1:carry-across-asset-classes:2}
Le taux d'intérêt domestique est de 1~\% et celui d'une devise étrangère de 5~\%. Calculer le carry d'une position longue sur le forward de cette devise, puis celui d'un
contrat à terme sur indice boursier avec un rendement de dividende de 2~\% lorsque le taux de financement vaut 4.5~\%~?
\end{exercise}

\begin{exercise}[$\star$]\label{exo:s1:carry-across-asset-classes:3}
Expliquer pourquoi la pente entre les deux premiers contrats est une mauvaise mesure du carry pour le gaz naturel.
\end{exercise}

\begin{exercise}[$\star\star$]\label{exo:s1:carry-across-asset-classes:4}
Le livre diversifié présente une volatilité de 10~\%, un ratio de Sharpe de 0.95 et perd 26.3~\% lors d'un krach
de cent jours. Calculer combien d'écarts-types à cent jours cette perte représente, et ce que le livre espérait gagner sur ces cent jours.
\end{exercise}

\begin{exercise}[$\star\star$]\label{exo:s1:carry-across-asset-classes:5}
Expliquer pourquoi le livre diversifié perd davantage que chacune des classes prises isolément lors du second krach.
\end{exercise}

\begin{exercise}[$\star\star$]\label{exo:s1:carry-across-asset-classes:6}
Expliquer, à l'aide de la décomposition de Koijen et coauteurs, pourquoi la part du carry effectivement perçu est «~toute la question~».
\end{exercise}

\begin{exercise}[$\star\star\star$]\label{exo:s1:carry-across-asset-classes:7}
\emph{Programmation.} Exécuter \texttt{books} avec le signal \texttt{near} et une prime de 1.0, qui lit le carry des matières premières à partir des deux premiers contrats.
Déterminer ce qu'il advient des ratios de Sharpe des matières premières et du livre diversifié, et pourquoi.
\end{exercise}

\begin{exercise}[$\star\star\star$]\label{exo:s1:carry-across-asset-classes:8}
\emph{Trouver la faille.} «~Notre livre de carry de change affiche un ratio de Sharpe de 1.0 depuis dix ans, avec un drawdown maximal de 8~\%~: c'est une stratégie à faible risque.~»
\end{exercise}

\section{Problème~: L'automne sur le marché synthétique}

\begin{problem}[{Problème du week-end --- le carry à travers un krach}]\label{pb:s1:carry-across-asset-classes:1}
L'univers synthétique du chapitre, la courbe du WTI et les archives publiques.

\medskip\noindent\textbf{Partie I --- Carry.}
\begin{enumerate}
\item Définir le carry d'une position sur contrat à terme et donner son nom dans chaque classe d'actifs.
\item Écrire la décomposition de l'espérance de rendement proposée par Koijen et coauteurs.
\item Définir une stratégie de carry et le \omterm{def:s1:carry-across-asset-classes:diversified}{carry diversifié}.
\item Expliquer comment le carry se lit sur les courbes synthétiques, et ce qui pose problème pour les matières premières saisonnières.
\end{enumerate}

\medskip\noindent\textbf{Partie II --- Les livres.}
\begin{enumerate}[resume]
\item Décrire les livres~: pondérations, coûts, dimensionnement.
\item Donner les ratios de Sharpe par classe et du livre diversifié pour $\lambda = 0.5$ et $\lambda = 1$.
\item Expliquer pourquoi le livre diversifié surpasse chaque classe.
\item Expliquer pourquoi le livre de change ne rapporte presque rien pour $\lambda = 0.5$.
\end{enumerate}

\medskip\noindent\textbf{Partie III --- Krachs.}
\begin{enumerate}[resume]
\item Définir un \omterm{def:s1:carry-across-asset-classes:crash}{krach du carry} et exposer ce que Brunnermeier, Nagel et Pedersen ont mis en évidence.
\item Exposer ce que Lustig, Roussanov et Verdelhan ont mis en évidence.
\item Donner les pertes du livre de change et du livre diversifié lors des trois krachs.
\item Expliquer ce que décrit l'encadré daté.
\end{enumerate}

\medskip\noindent\textbf{Partie IV --- Le verdict.}
\begin{enumerate}[resume]
\item Énoncer le \emph{résultat nommé}~: le ratio de Sharpe du livre de \omterm{def:s1:carry-across-asset-classes:diversified}{carry diversifié} et sa perte lors du krach, comparés à ceux de chaque classe prise isolément.
\item Expliquer comment le carry et le trend se combinent.
\item Donner le carry du WTI de 1985 à 2024, et dire ce que la règle du signe a rapporté.
\item Dire à quoi correspondent les extrêmes du carry du WTI.
\item Expliquer comment réaliser un backtest honnête d'une stratégie de carry.
\item Dire quelle fiche de stratégie est la plus exposée aux limites de stockage.
\item Exposer ce que Koijen et coauteurs ont mis en évidence sur le carry à travers les classes d'actifs en récession.
\item Dire ce pour quoi un investisseur en carry est rémunéré, en une phrase.
\end{enumerate}
\end{problem}

\section{Questions d'entretien}

\begin{interviewq}[$\star$ \iqroles{chercheur}]\label{iq:s1:carry-across-asset-classes:1}
Qu'est-ce que le carry pour une devise, un contrat à terme sur obligations et un contrat à terme sur matière première~?
\end{interviewq}

\begin{interviewq}[$\star\star$ \iqroles{chercheur}]\label{iq:s1:carry-across-asset-classes:2}
Pourquoi la parité des taux d'intérêt non couverte impliquerait-elle que le carry de change ne rapporte rien, et que disent les données empiriques~?
\end{interviewq}

\begin{interviewq}[$\star\star$ \iqroles{trader}]\label{iq:s1:carry-across-asset-classes:3}
La volatilité bondit et le yen gagne 3~\% en une journée. Que faites-vous de votre livre de carry de change~?
\end{interviewq}

\begin{interviewq}[$\star\star$ \iqroles{risque}]\label{iq:s1:carry-across-asset-classes:4}
Comment mesureriez-vous le risque de krach d'un livre de carry qui n'a jamais connu de mauvais mois dans son histoire~?
\end{interviewq}

\begin{interviewq}[$\star\star$ \iqroles{développeur}]\label{iq:s1:carry-across-asset-classes:5}
De quelles données un livre de \omterm{def:s1:carry-across-asset-classes:diversified}{carry diversifié} a-t-il besoin chaque jour, et que peut-il mal tourner~?
\end{interviewq}

\begin{interviewq}[$\star\star\star$ \iqroles{chercheur}]\label{iq:s1:carry-across-asset-classes:6}
Le cours d'un contrat à terme vaut $F = S e^{-c\tau}$ avec un carry $c$ constant et un spot $S$ suivant une martingale. Montrer que le rendement espéré de la position vaut $c$
par unité de temps, puis déterminer le rendement espéré lorsque, à la place, $E[dS/S] = -(1 - \lambda) c\, dt$.
\end{interviewq}
